\chapter{Exact Normalization and Quantization}\label{ch08:quantization}

\section{The problem a quantizer inherits}\label{ch08:sec:problem}

Of the eight conflicts preserved in the source audit, the quantization
conflict is the smallest in surface area and the most instructive in kind.
The Quantification Guide displays four score bands with two-decimal
endpoints --- $0.00$--$2.49$, $2.50$--$4.99$, $5.00$--$7.49$,
$7.50$--$10.00$ --- and provides no rule for values of higher precision.
The conflict is not between two sources; it is between the display and the
mathematics. Two-decimal endpoint pairs like $2.49$ and $2.50$ presuppose
that scores arrive already rounded to two decimals, and are silent about
every value in the open gap between them. A value of $2.495$, a value of
$2.4999$, a value produced by dividing two exact decimals whose quotient
does not terminate at two places --- all fall outside the displayed
partition. A classification system whose bands are defined on displayed
values inherits the rounding rule of its display layer as an invisible
classification parameter, and different rounding conventions then classify
the same underlying value differently.

The candidate's resolution is structural: classification is computed from
the exact pre-rounding value, over half-open bins whose cut points are the
exact rationals $5/2$, $5$, and $15/2$; the displayed value is derived
afterward and never re-enters classification. This chapter states the exact
formulas as implemented, proves the boundary-preservation property they
were designed for, works the boundary cases numerically, and explains
precisely why the Guide's two-decimal bands underdetermine classification.
It closes with the trend-tag discipline, which extends the same
exactness-first posture from single values to comparisons across values.

As everywhere in this monograph, the status frame governs: the quantizer is
tested software operating on typed records
(\textsc{open / ring 4 / not licensed / held}); no normalized score
reported here derives from an admitted live-model measurement, and all
numeric values in the worked examples are illustrative constructions.

\section{Exact normalization as implemented}\label{ch08:sec:exact}

\begin{definition}[Normalization formulas]\label{ch08:def:formulas}
Let $\mathrm{raw} \ge 0$ be an instrumented flow magnitude in
\textsc{activation} units at a registered test point, and let
$I_{\mathrm{FS}} > 0$ be the calibration profile's full-scale value for
that test point, produced under the profile's declared full-scale method.
The implemented normalization computes three quantities:
\[
q_{10}^{\mathrm{exact}} \;=\; 10 \cdot \frac{\mathrm{raw}}{I_{\mathrm{FS}}},
\qquad
\mathrm{overload} \;=\; \frac{\mathrm{raw}}{I_{\mathrm{FS}}},
\qquad
q_{10}^{\mathrm{disp}} \;=\; \min\bigl(q_{10}^{\mathrm{exact}},\, 10\bigr).
\]
The displayed score $q_{10}^{\mathrm{disp}}$ is clipped at $10$; the
overload ratio is never clipped; and the exact value
$q_{10}^{\mathrm{exact}}$ is retained as an exact ratio object --- numerator
$\mathrm{raw}$, denominator $I_{\mathrm{FS}}$, scale $10$ --- whose value is
evaluated in exact rational arithmetic, not in binary floating point.
\end{definition}

Three implementation details make the definition load-bearing rather than
decorative.

First, the three quantities are constructed together, inside a factory, and
sealed into a single result object (\texttt{NormalizedFlow}) that cannot be
built by any other path: the constructor validates a factory token and
rejects direct construction. The result object re-verifies, at admission,
the arithmetic identity between its parts --- that the stored displayed
value equals $\min(10, 10 \cdot \mathrm{raw}/I_{\mathrm{FS}})$ and the
stored overload ratio equals $\mathrm{raw}/I_{\mathrm{FS}}$, both recomputed
from the exact ratio --- so a normalized record whose display has drifted
from its exact value is unrepresentable, not merely unlikely.

Second, the inputs are typed evidence, not numbers. The raw flow must be a
\texttt{FLOW\_MAGNITUDE} in \textsc{activation} units, of measurement kind
\texttt{INSTRUMENTED}, with a recorded sample size, bound to the same test
point for which the profile's full-scale value is registered; the
full-scale value itself is resolved from the profile and re-verified against
the evidence registry. A quantity of the wrong kind, unit, provenance, or
test point is rejected before any arithmetic occurs. The mixed-scale error
preserved in the source audit --- the Guide's $\mathrm{sat}=8.5$ read
against a $0$--$1$ fraction convention --- is in this sense not merely
avoided but unrepresentable: a Q10 score, a Q100 constraint score, a
fraction, and an activation magnitude are four different unit kinds, and
the unit-admissibility table refuses the cross-assignments.

Third, the exact value is carried as a ratio of decimals evaluated through
exact rational arithmetic (numerator and denominator lifted to arbitrary-precision
rationals), so that the band comparisons of the next section are performed
on a representation in which equality and order at the cut points are exact
questions with exact answers. No binary float participates in
classification.

\begin{remark}[Why the overload ratio is a separate quantity]\label{ch08:rem:overload}
Clipping the display at $10$ is a presentation decision inherited from the
Q10 score convention; it deliberately destroys information above full
scale. The overload ratio preserves exactly the destroyed information: a
reading of $\mathrm{overload} = 1.2$ records that the raw flow exceeded the
declared full scale by twenty percent, which the clipped display
$q_{10}^{\mathrm{disp}} = 10$ cannot distinguish from a reading exactly at
full scale. Because the full-scale value is itself method-relative --- the
profile's declared \texttt{IFSMethod} is one of three non-interchangeable
procedures, per the preserved full-scale conflict --- an overload above $1$
is ambiguous between a genuinely extreme response and a full-scale value
selected too low for the point; the ratio does not resolve the ambiguity,
it preserves it in a typed, dimensionless quantity
(\texttt{OVERLOAD\_RATIO}) that downstream analysis can interrogate.
Overload is also the natural bridge to the saturation discussions of
Chapter~\ref{ch07:classification}: a saturated region label with overload
near $1$ reads differently from the same label with overload far above $1$,
and neither reading is licensed automatically.
\end{remark}

\section{Representation: decimals in, rationals underneath}\label{ch08:sec:representation}

The exactness claims of this chapter rest on a representation choice
that the candidate makes globally and the quantizer inherits: every
numeric value in the evidence layer is an exact decimal --- a scaled
integer, not a binary float --- and every ratio that classification
consumes is carried as an explicit numerator-denominator pair evaluated
in arbitrary-precision rational arithmetic. The choice forecloses three
familiar failure classes at once.

First, \emph{representation error at the boundary}. The cut point $2.5$
is exactly representable in decimal and in rational arithmetic; a value
like $2.4999$ is likewise exact. In binary floating point, neither the
inputs nor the cut points of a decimal band scheme are generally
representable, and a comparison at the boundary inherits up to half an
ulp of silent displacement --- small, but structurally identical to the
rendering gap of Section~\ref{ch08:sec:underdetermine}, and
platform-honest rather than convention-honest. The rational route makes
the comparison a cross-multiplication of integers, which is exact on
every platform.

Second, \emph{division before comparison}. A naive implementation would
compute $10 \cdot \mathrm{raw} / I_{\mathrm{FS}}$ as a quotient first
and compare afterward, forcing a precision decision (how many digits of
quotient?) before the comparison that needed the answer. The
implemented route stores the ratio unevaluated --- numerator, denominator,
scale --- and compares by cross-multiplication, so no quotient precision
is ever chosen for classification purposes. The displayed quotient, by
contrast, is computed in the candidate's declared decimal arithmetic
context, where a precision \emph{is} chosen; the asymmetry is the whole
design: presentation gets a precision, classification gets none because
it needs none.

Third, \emph{drift between stored value and stored band}. Because the
band is a function, not a field --- it is recomputed from the exact
ratio on demand, and the sealed result object verifies its internal
identities at construction --- there is no stored band to fall out of
sync with a stored value. Nothing in the evidence model can hold a
record asserting exact value $2.4$ and band $Q2$; the assertion is not
storable, and this unstorability is the quantizer's contribution to the
candidate's general no-bare-numbers discipline.

\section{Half-open bins}\label{ch08:sec:bins}

\begin{definition}[Q4 bands]\label{ch08:def:bands}
Band assignment is a total function of the exact value
$q_{10}^{\mathrm{exact}}$, through half-open intervals closed on the left:
\[
[0, 2.5) \mapsto Q1, \qquad
[2.5, 5) \mapsto Q2, \qquad
[5, 7.5) \mapsto Q3, \qquad
[7.5, \infty) \mapsto Q4 .
\]
Equivalently, as implemented: the value is compared in exact rational
arithmetic against $5/2$, then $5$, then $15/2$, returning the band at the
first strict-inequality success and $Q4$ otherwise.
\end{definition}

Three observations. First, the bins partition the entire non-negative ray:
every admissible exact value, including values above $10$ produced by
overload, receives exactly one band. The implementation's top band is
$[7.5, \infty)$, not $[7.5, 10]$: an overloaded value of $12$ is a $Q4$
value with an overload ratio recorded alongside, rather than an error or an
out-of-range special case. (The audit's prose summary states the top bin as
$[7.5, 10]$; the implemented comparison function has no upper test, and the
two agree on all clipped displays while differing on the exact values above
$10$ that only the overload path produces. This monograph records the
implemented behavior.) Second, each cut point belongs to the upper band:
$2.5$ is $Q2$, $5$ is $Q3$, $7.5$ is $Q4$. The choice of upward assignment
is a convention, but having \emph{some} exact convention is not: it is what
makes band assignment a function. Third, the cut points are exact rationals,
so membership questions are decidable without tolerance: there is no
epsilon, no rounding mode, and no platform dependence in the band function.

\section{The boundary-preservation proposition}\label{ch08:sec:boundary}

\begin{proposition}[Boundary preservation]\label{ch08:prop:boundary}
Band assignment is computed from $q_{10}^{\mathrm{exact}}$ by exact
comparison against $2.5$, $5$, and $7.5$. Display derivation --- clipping at
$10$ and any subsequent rounding for presentation --- cannot move a value
across a band boundary, because the displayed value never re-enters
classification.
\end{proposition}

\begin{designargument}
The claim decomposes into three parts.

(i) \emph{The band function reads only the exact value.} By construction,
the implemented band function accepts only the sealed normalized-flow
result and evaluates only its exact-ratio component; the displayed value
and the overload ratio are not arguments to any comparison. The band of a
record is therefore a function of
$(\mathrm{raw}, I_{\mathrm{FS}})$ alone, through the exact quotient.

(ii) \emph{Exact comparison is total and unambiguous at the cut points.}
The exact value is a non-negative rational with an explicit numerator and
denominator; the cut points $5/2$, $5$, and $15/2$ are rationals. For any
two rationals, exactly one of $<$, $=$, $>$ holds, and the comparison is
computed exactly in integer arithmetic on cross-multiplied numerators and
denominators. Hence for every representable input the guarded cascade
($< 5/2$, else $< 5$, else $< 15/2$, else $Q4$) selects exactly one band.
In particular, an input exactly equal to a cut point fails the strict
inequality at that cut point and is assigned to the upper band; equality is
not a boundary anomaly but an ordinary, deterministically handled case.

(iii) \emph{Display is derived afterward and is inert.} The displayed value
$q_{10}^{\mathrm{disp}} = \min(q_{10}^{\mathrm{exact}}, 10)$, and any
two-decimal (or other) rounding applied for presentation, are computed from
the exact value after classification and are consumed only by presentation
surfaces. Since no classification path accepts a displayed value as input
--- the band function's type signature admits only the sealed result, whose
exact component is verified at admission against its own display --- there
is no composition of implemented functions in which a rounded value
determines a band. A display rounding that crosses a cut point (say,
$2.4999 \rightarrow 2.50$) therefore changes what a human reads, and
nothing else.

Combining (i)--(iii): the composite map from $(\mathrm{raw},
I_{\mathrm{FS}})$ to (band, display) computes the band from the exact
quotient by total exact comparison and computes the display independently
downstream; no display operation participates in, or can alter, band
assignment.
\end{designargument}

Two small consequences are worth separating from the proposition,
because they answer questions a careful reader will ask next.

\begin{lemma}[Display is monotone and non-injective]\label{ch08:lem:display}
The display map $q_{10}^{\mathrm{exact}} \mapsto q_{10}^{\mathrm{disp}}
= \min(q_{10}^{\mathrm{exact}}, 10)$ is order-preserving (weakly) and
not injective: it is the identity on $[0, 10]$ and constant on
$[10, \infty)$. Consequently the pair (display, overload) determines
the exact value --- the overload ratio is injective on the whole ray
--- while the display alone determines it only below full scale.
\end{lemma}

\begin{proof}
Monotonicity and the two-piece description are immediate from the
definition of $\min$. Injectivity of the overload ratio holds because
$\mathrm{overload} = q_{10}^{\mathrm{exact}} / 10$ identically ---
overload is the same exact quotient at unit scale --- so the exact
value is recoverable from it by multiplication by $10$; the display
coincides with the exact value on $[0,10]$ and discards it above.
\end{proof}

\begin{corollary}[Band recoverability]\label{ch08:cor:recover}
The band of a record is recoverable from its overload ratio alone
(compare against $1/4$, $1/2$, $3/4$), and from its display alone if
and only if the display is strictly below $10$ or the band is $Q4$.
In particular, a consumer holding only clipped displays can still
classify correctly, because every clipped record displays $10$ and
$10 \ge 7.5$ lies in $Q4$; what such a consumer cannot do is anything
finer than banding in the overload regime.
\end{corollary}

Table~\ref{ch08:tab:bands} collects the band structure with its
boundary behavior in one place.

\begin{table}[htbp]
\centering
\begin{tabular}{@{}lllll@{}}
\toprule
Band & Exact interval & Cut point at left & Guide display band & Gap under two decimals \\
\midrule
$Q1$ & $[0, 2.5)$ & $0$ (closed) & $0.00$--$2.49$ & $(2.49, 2.50)$ unruled \\
$Q2$ & $[2.5, 5)$ & $2.5 \in Q2$ & $2.50$--$4.99$ & $(4.99, 5.00)$ unruled \\
$Q3$ & $[5, 7.5)$ & $5 \in Q3$ & $5.00$--$7.49$ & $(7.49, 7.50)$ unruled \\
$Q4$ & $[7.5, \infty)$ & $7.5 \in Q4$ & $7.50$--$10.00$ & overload collapsed at $10.00$ \\
\bottomrule
\end{tabular}
\caption{Exact half-open bands versus the Guide's printed two-decimal
display bands. Every printed endpoint pair leaves an unruled open gap;
the implemented cut points rule every value; the top band differs
additionally in its treatment of overload.}
\label{ch08:tab:bands}
\end{table}

\begin{remark}\label{ch08:rem:proof-scope}
The proposition is a software-structural statement, verified by the
candidate's test suite at the level of golden boundary fixtures; it is not
a metrological claim. It says that the quantizer cannot \emph{add} a
boundary error to whatever measurement error the raw value already
carries; it does not say the raw value is right. Uncertainty attached to
the raw flow propagates through normalization as interval and standard
uncertainty on the ratio (Section~\ref{ch08:sec:uncertainty}); a raw value
whose uncertainty interval straddles a cut point yields a band assignment
that is exact for the point value and undecided across the interval, and
honest reporting of that situation is the trend-tag discipline's concern
(Section~\ref{ch08:sec:trend}).
\end{remark}

\section{Worked examples}\label{ch08:sec:worked}

All values below are illustrative constructions, labeled as such; none is a
measurement.

\begin{example}[Overload: $\mathrm{raw}/I_{\mathrm{FS}} = 1.2$]\label{ch08:ex:overload}
Take $\mathrm{raw} = 6$ activation units at a test point whose declared
full scale is $I_{\mathrm{FS}} = 5$. Then
\[
q_{10}^{\mathrm{exact}} = 10 \cdot \tfrac{6}{5} = 12, \qquad
\mathrm{overload} = \tfrac{6}{5} = 1.2, \qquad
q_{10}^{\mathrm{disp}} = \min(12, 10) = 10 .
\]
Band assignment reads the exact value $12$: since $12 \ge 15/2$, the band
is $Q4$. The record thus carries four coordinated facts --- exact $12$,
display $10$, overload $1.2$, band $Q4$ --- and the sealed result object
verifies their mutual consistency at construction. Note what the display
alone would have lost: a reader of ``$10.00$'' cannot distinguish this
record from a reading exactly at full scale, but the record itself can,
and any consumer that reads the typed object rather than its rendering
retains the distinction.
\end{example}

\begin{example}[Exact cut point: $q_{10}^{\mathrm{exact}} = 2.5$]\label{ch08:ex:cut}
Take $\mathrm{raw} = 1$ and $I_{\mathrm{FS}} = 4$, so
$q_{10}^{\mathrm{exact}} = 10/4 = 5/2$ exactly. The first guard
($< 5/2$) fails on equality; the second guard ($< 5$) succeeds; the band
is $Q2$. The half-open convention assigns every cut point upward, so
$2.5$ is $Q2$, $5$ is $Q3$, and $7.5$ is $Q4$, deterministically. Under
the Guide's displayed bands the value $2.50$ also reads as the second band
($2.50$--$4.99$), so at exactly representable two-decimal values the
implemented convention and the Guide's display agree; the divergence
appears strictly between the displayed endpoints, as the next example
shows.
\end{example}

\begin{example}[Near-boundary value: $q_{10}^{\mathrm{exact}} = 2.4999$]\label{ch08:ex:near}
Take $\mathrm{raw} = 2.4999$ and $I_{\mathrm{FS}} = 10$, so
$q_{10}^{\mathrm{exact}} = 2.4999$ exactly. Since $2.4999 < 5/2$, the band
is $Q1$. A two-decimal display that rounds half up renders this value as
$2.50$ --- the printed lower endpoint of the Guide's \emph{second} band.
A reader classifying from the display would place the record in $Q2$; the
implemented classifier, reading the exact value, places it in $Q1$; and by
Proposition~\ref{ch08:prop:boundary} the rendered ``$2.50$'' has no way to
leak back into classification. The example generalizes: every value in
$[2.495, 2.5)$ rounds to $2.50$ at two decimals while lying exactly in
$Q1$, so the display-classification gap is an interval of positive width
at every cut point, not an isolated pathological value.
\end{example}

\begin{example}[Gap value under the Guide's bands: $2.495$]\label{ch08:ex:gap}
The value $2.495$ lies in the open gap between the Guide's printed
endpoints $2.49$ and $2.50$: it belongs to neither $0.00$--$2.49$ nor
$2.50$--$4.99$ as literally written. The Guide provides no rule for it.
The implemented convention answers instantly: $2.495 < 5/2$, band $Q1$.
The point of the example is not that the Guide's authors intended
otherwise --- plausibly they intended truncation, or rounding, or never
contemplated the gap --- but that ``plausibly intended'' is not a
classification rule, and three incompatible completions (truncate at two
decimals, round half up, round half even) each classify parts of the gap
differently.
\end{example}

\begin{example}[Remaining cut points and degenerate inputs]\label{ch08:ex:more}
All values illustrative. (i) $\mathrm{raw} = 2$, $I_{\mathrm{FS}} = 4$:
exact value $5$ exactly; the first two guards fail (at the second, on
equality); band $Q3$. (ii) $\mathrm{raw} = 7.4999$,
$I_{\mathrm{FS}} = 10$: exact value $7.4999 < 15/2$; band $Q3$,
although a two-decimal display renders $7.50$, the printed floor of
the Guide's top band --- the $Q4$-boundary instance of the rendering
gap. (iii) $\mathrm{raw} = 0$: exact value $0$; band $Q1$ by the closed
left end of the first bin; overload $0$; and note that a zero
\emph{observed} flow is banded without difficulty even though ratio
diagnostics elsewhere in the pipeline would censor at a zero-compatible
observation --- banding divides by the full scale, which is strictly
positive by profile construction, never by the observation.
(iv) $\mathrm{raw} = 1.874975$, $I_{\mathrm{FS}} = 7.4999$: exact value
$10 \cdot 1.874975 / 7.4999 = 2.5$ exactly (the numerator is
$0.25 \times 7.4999$); band $Q2$ by upward assignment --- an example
chosen to show that cut-point equality arises from unremarkable-looking
inputs, not only from round numbers.
\end{example}

\begin{example}[Uncertainty straddling a cut point]\label{ch08:ex:straddle}
All values illustrative. A raw flow of $1.24$ with standard
uncertainty $0.03$ against $I_{\mathrm{FS}} = 5$ (with standard
uncertainty $0.25$) yields exact value $2.48$ and, after propagation, an
interval on the Q10 coordinate of $[2.34, 2.62]$ (first-order
combination: $2.48\cdot\sqrt{(0.03/1.24)^2 + (0.25/5)^2} \approx 0.14$). The point
value is $Q1$; the interval spans the $2.5$ cut. The record reports
all three facts. What the record does not do is report ``$Q1$--$Q2$''
or a probability over bands: banding stays exact on the point value,
and the interval stays available to any consumer whose question is
distributional. A trend or comparison built on this record must use
the interval under the declared uncertainty rule of
Section~\ref{ch08:sec:trend}; a display built on it will render
$2.48$; and neither operation can move the recorded band.
\end{example}

\section{Why two-decimal display bands underdetermine classification}\label{ch08:sec:underdetermine}

The preceding examples can now be assembled into the general argument,
which is the content of the preserved quantization conflict.

A band scheme defined on two-decimal displays is a partition of the
\emph{rendered} values $\{0.00, 0.01, \ldots\}$, not of the underlying
value space. To classify an underlying value, the scheme must be composed
with a rendering rule, and the composition --- not the printed bands ---
is the actual classifier. The Guide prints the bands and omits the
rendering rule; therefore the Guide's scheme is not one classifier but a
family of classifiers indexed by the omitted rule, and members of the
family disagree on every interval of the form
$[c - 0.005, c)$ for each cut point $c \in \{2.5, 5.0, 7.5\}$ (under
round-half-up rendering; the disagreeing intervals shift with the
rendering convention but never vanish). The underdetermination is thus
structural, not editorial: no choice of printed two-decimal endpoints can
close it, because the gap lives between the printed values.

The conflict is sharpest at the top band --- the Q4-boundaries conflict of
the source audit --- for two compounding reasons. First, the Guide's top
band is printed as $7.50$--$10.00$, closed above, which presupposes that
no value exceeds $10.00$; but the normalization that produces the score
admits raw values above full scale, and the clipped display maps all of
them to exactly $10.00$. Under a display-defined scheme, the entire
overload regime collapses into the top endpoint, and the scheme cannot
even express the question ``by how much was full scale exceeded.'' The
implemented scheme keeps the top band half-open to infinity on the exact
value and records overload separately, so the collapse happens only in
presentation. Second, the boundary $7.50$ itself inherits the rendering
gap like every other cut point: a value of $7.4999$ renders as $7.50$ and
would read as top-band under the display while lying exactly in $Q3$. A
two-decimal scheme therefore misclassifies (relative to the exact
convention) in both directions at the Q4 boundary --- upward across the
cut by rounding, and degenerately at the top by clipping --- and these are
precisely the two failure modes the exact-value-plus-overload design
eliminates.

\begin{finding}[Display bands are not a classifier]\label{ch08:find:display}
The Guide's two-decimal bands, as printed, do not determine a
classification function on exact values; they determine one only relative
to an unstated rendering rule, and distinct reasonable rules yield
distinct classifications on intervals of positive width at every band
boundary. The implemented candidate closes the gap by fiat --- exact
values, half-open bins, upward cut-point assignment, display derived
afterward --- and Proposition~\ref{ch08:prop:boundary} certifies that the
closure is complete: no rendering choice can re-open it.
\end{finding}

It bears stating why this matters at monograph length rather than as a
footnote. Quantized bands are the layer at which numeric evidence becomes
categorical language --- ``$Q1$'' reads as \emph{low} and ``$Q4$'' as
\emph{high} --- and categorical language is what travels into summaries,
dashboards, and decisions. A classification gap at the band boundary is
therefore not a cosmetic defect: it is a mechanism by which two honest
readers of the same underlying value can carry incompatible categorical
claims forward, each traceable to the same source, with the disagreement
invisible unless the rendering rule is recovered. The dose-response
tradition's practice of modeling the continuous response directly
\cite{ritz} --- together with the broader statistical caution against
premature categorization of continuous quantities --- is the same
discipline in statistical form;
the candidate implements it as a type-level separation between an exact
ratio that classifies and a display that does not.

\section{Uncertainty through normalization}\label{ch08:sec:uncertainty}

Normalization does not launder uncertainty. The implemented path requires
uncertainty on the raw flow and on the full-scale value, and propagates
both through the ratio: interval bounds by the monotone quotient rule
(lower display bound from lower numerator over upper denominator, upper
from upper over lower), standard uncertainty by the first-order
combination of the two relative contributions, and clipping applied to the
Q10 display bounds at $10$ while the overload ratio's bounds remain
unclipped. The confidence level attached to the normalized quantities is
the minimum of the input confidence levels. That label is an upper bound
on joint coverage, not a conservative floor: the quotient interval covers
the true ratio only when both input intervals cover, an event with
probability at most the minimum (approximately the product under
independence), so reading the recorded level as attained coverage would
overstate the evidence; replacing the convention with a Bonferroni-style
lower bound $\max(0,\,1-\sum_i(1-c_i))$ is an open revision item.
Treating the normalized best value as if it were the truth,
with the uncertainty discarded at the seam, is the calibration sin the
wider candidate architecture is built to refuse \cite{kennedyohagan}; the
quantizer participates by never emitting a normalized value object without
its propagated uncertainty attached.

Propagated uncertainty and exact banding answer different
questions and are easily conflated. Banding is exact on the point value:
a record whose exact value is $2.4999$ is a $Q1$ record, full stop. But if
that record's propagated interval spans $[2.3, 2.7]$, then the statement
``the underlying quantity is in $Q1$'' is supported only at the point
estimate, not across the interval, and the record carries everything a
consumer needs to see this: exact value, band, and interval, jointly. The
candidate's position is that the band function must not blur itself to
accommodate uncertainty (a fuzzy band is a second, worse rendering rule);
instead, uncertainty-aware statements are pushed to the comparison layer,
where they are made explicit --- which is the trend-tag rule.

\section{The quantizer and the mixed-scale conflict}\label{ch08:sec:mixedscale}

The quantization conflict is the audit's eighth; its first --- the
mixed-scale conflict --- is enforced partly here, and the connection
is stated explicitly here because the two conflicts compound. The
Guide's worked example uses $\mathrm{sat} = 8.5$ as a component
threshold on a page whose conventions make saturation a $0$--$1$
fraction; Appendix~I identifies the value as a mixed-scale error whose
intended magnitude is $0.85$. On a Q10 coordinate, $8.5$ and $0.85$
fall in different bands ($Q4$ and $Q1$); on a fraction coordinate,
$8.5$ is not even admissible. A quantizer that accepted bare numbers
would classify the error fluently --- $8.5$ is a perfectly bandable
float --- and the misclassification would propagate wearing a band
label of unimpeachable exactness.

The implemented defense is the unit ontology upstream of the band
function. The quantizer's input is a typed measurement whose quantity
kind (\texttt{FLOW\_MAGNITUDE}) and unit (\textsc{activation}) are
checked before normalization; its output Q10 score is a
\texttt{NORMALIZED\_FLOW} in unit \texttt{Q10}, range-checked to
$[0, 10]$; fractions are a distinct unit kind range-checked to
$[0, 1]$; Q100 constraint scores are a third, range-checked to
$[0, 100]$; and the unit-admissibility table permits each quantity kind
only its declared units, with no conversion function anywhere in the
layer. A saturation fraction of $8.5$ is a construction-time validation
error; a Q10 score used where a fraction is required is a type error;
and the Guide's $\mathrm{sat}=8.5$ is therefore not a value the
candidate handles \emph{incorrectly} but a value it cannot admit at
all. Exact banding sits behind these gates: exactness is only worth
proving about values whose kind and scale are already beyond dispute,
which is why this chapter's guarantees are stated for typed inputs and
why the mixed-scale defense, though implemented elsewhere, is properly
part of the quantization story.

\section{Trend tags}\label{ch08:sec:trend}

A single normalized value invites a band word; a pair of them invites a
trend word --- \emph{rising}, \emph{falling}, \emph{stable} --- and trend
words smuggle in even more unstated structure than band words do. A trend
is a claim about two or more values \emph{of the same quantity, on the
same axis, over a stated interval}, and it is empty or misleading if any
of those qualifiers is left implicit: values from different test points do
not form a trend; values under different profiles do not form a trend;
a difference smaller than the propagated uncertainty does not evidence a
trend; and a difference smaller than a declared minimum effect is not a
trend worth naming even if statistically resolvable.

\begin{ruleenv}[Trend-tag admission]\label{ch08:rule:trend}
A trend tag is admissible only when it declares, explicitly: (i) the axis
--- which typed quantity, at which test point, under which profile digest,
the profile identity fixing model, method, and full-scale convention;
(ii) the interval --- the ordered replay frame over which the compared
values were produced, with the caller-supplied replay semantics of
Chapter~\ref{ch07:classification} understood; (iii) the minimum effect ---
a declared threshold below which differences are reported as no trend,
fixed before comparison; and (iv) the uncertainty rule --- how the
propagated uncertainties of the compared values enter the judgment (for
example, non-overlap of intervals at the stated confidence, or a stated
multiple of combined standard uncertainty). A tag missing any component is
not a degraded trend statement; it is inadmissible.
\end{ruleenv}

The rule is the comparison-layer twin of the banding discipline: just as a
band must come from the exact value and not the rendering, a trend must
come from a declared comparison rule and not from the visual suggestion of
two displayed numbers. The two-decimal display is again the cautionary
instrument: displays of $2.50$ and $2.50$ can cover exact values $2.4950$
and $2.5049$ --- a difference of $0.0099$ in the exact coordinate rendered
as perfect stability --- while displays of $2.49$ and $2.50$ can cover
exact values differing by one part in $10^4$, rendered as a visible step.
Display-level trend reading therefore manufactures both false stability
and false movement at the rendering quantum, and the admission rule
forecloses it by requiring the comparison to be defined on the exact
values with an explicit uncertainty rule.

\begin{example}[An admissible trend tag, fully declared]\label{ch08:ex:trend}
All values illustrative. An admissible tag reads, in full: ``Axis:
normalized flow (\texttt{I\_Q}, Q10) at test point
\texttt{tp-}$\langle$digest$\rangle$ under profile
$\langle$digest$\rangle$. Interval: replay frames $t_1 < t_2$, both
inside the profile validity window. Minimum effect: $0.5$ Q10 units,
declared before comparison. Uncertainty rule: report \emph{rising} only
if the lower bound of the later value exceeds the upper bound of the
earlier value and the difference of point values exceeds the minimum
effect.'' Under this tag, exact values $6.1 \to 6.9$ with intervals
$[5.9, 6.3]$ and $[6.7, 7.1]$ read \emph{rising} (intervals disjoint,
difference $0.8 > 0.5$); exact values $6.1 \to 6.5$ with the same
half-widths read \emph{no trend} (difference below minimum effect,
whatever the intervals); and exact values $6.1 \to 6.9$ with intervals
$[5.5, 6.7]$ and $[6.3, 7.5]$ read \emph{no trend} (intervals overlap),
even though the displayed values suggest movement. Delete any one of
the four declarations and each verdict becomes contestable, which is
the rule's point.
\end{example}

No trend tag over live-model data is asserted anywhere in this monograph;
the rule is stated for the discipline it imposes on future admitted data,
and its components (axis, interval, minimum effect, uncertainty rule) are
deliberately the same components a preregistered comparison would freeze
in advance. The kinship with dose-response practice is direct: a trend
over dose is a model claim, and the modeling tradition fits and tests it
on the continuous response with declared error structure \cite{ritz},
rather than reading it off category transitions; the trend-tag
rule is that insistence, restated for the replay-frame axis.

\section{What the quantization layer does and does not establish}\label{ch08:sec:close}

\begin{boundary}[Scope of the exactness guarantees]\label{ch08:bound:scope}
The guarantees of this chapter are guarantees about representation and
classification mechanics: exact normalization, total and deterministic
banding, display inertness, uncertainty propagation, and comparison
admission. They are not guarantees about measurement validity. An exact
$Q4$ band over a mismeasured raw flow is an exact classification of a
wrong number; a faithfully propagated uncertainty is only as honest as
the input uncertainty model; and the full-scale denominators remain
method-relative under the preserved full-scale conflict, so cross-profile
band comparisons are undefined absent a stated comparison rule. The
quantization layer's contribution is negative and structural: of the ways
a score pipeline can silently corrupt a classification --- rounding across
boundaries, clipping away overload, unit slippage, display feedback ---
none is available here, and the tests that pin each closure are part of
the bounded software evidence at the Gate~4 receipt.
\end{boundary}
